\documentclass[11pt, a4paper]{article}

% ─── Packages ──────────────────────────────────────────────────────────────────
\usepackage[utf8]{inputenc}
\usepackage[T1]{fontenc}
\usepackage{amsmath, amssymb, amsthm, mathtools}
\usepackage{bm}                         % Bold math symbols
\usepackage{bbm}                        % Blackboard bold for indicators
\usepackage{physics}                    % Derivatives, norms, brackets
\usepackage{algorithm, algpseudocode}   % Pseudocode
\usepackage{booktabs}                   % Tables
\usepackage{graphicx}                   % Figures
\usepackage{hyperref}                   % Clickable references
\usepackage[margin=2.5cm]{geometry}
\usepackage{enumitem}
\usepackage{xcolor}
\usepackage{tcolorbox}                  % Colored boxes for key results
\usepackage{listings}                   % Code snippets
\usepackage[backend=biber, style=authoryear, sorting=nyt]{biblatex}

% ─── Theorem environments ──────────────────────────────────────────────────────
\theoremstyle{definition}
\newtheorem{definition}{Definition}[section]
\newtheorem{example}{Example}[section]
\theoremstyle{plain}
\newtheorem{theorem}{Theorem}[section]
\newtheorem{proposition}{Proposition}[section]
\newtheorem{lemma}{Lemma}[section]
\newtheorem{corollary}{Corollary}[section]
\theoremstyle{remark}
\newtheorem{remark}{Remark}[section]

% ─── Key result box ────────────────────────────────────────────────────────────
\newtcolorbox{keyresult}[1][]{
    colback=blue!5!white,
    colframe=blue!60!black,
    fonttitle=\bfseries,
    title={Key Result},
    #1
}

% ─── Code-link box (connects math to implementation) ──────────────────────────
\newtcolorbox{codelink}[1][]{
    colback=green!5!white,
    colframe=green!50!black,
    fonttitle=\bfseries,
    title={Implementation Reference},
    #1
}

% ─── Code listing style ───────────────────────────────────────────────────────
\lstset{
    language=Python,
    basicstyle=\ttfamily\small,
    keywordstyle=\color{blue},
    commentstyle=\color{gray},
    stringstyle=\color{red!70!black},
    numbers=left,
    numberstyle=\tiny\color{gray},
    frame=single,
    breaklines=true,
}

% ─── Custom commands ───────────────────────────────────────────────────────────
\newcommand{\VaR}{\operatorname{VaR}}
\newcommand{\ES}{\operatorname{ES}}
\newcommand{\CVaR}{\operatorname{CVaR}}
\newcommand{\EV}{\mathbb{E}}           % Expectation
\newcommand{\PV}{\mathbb{P}}           % Probability
\newcommand{\QV}{\mathbb{Q}}           % Risk-neutral measure
\newcommand{\RV}{\mathbb{R}}           % Real numbers
\newcommand{\ind}{\mathbbm{1}}         % Indicator function
\newcommand{\cov}{\operatorname{Cov}}
\newcommand{\corr}{\operatorname{Corr}}
\newcommand{\var}{\operatorname{Var}}

% ─── Document ──────────────────────────────────────────────────────────────────
% TEMPLATE: Copy this file to your project's docs/ directory and customize.
% Replace the placeholders below.

\title{%
    \textbf{Project XX: [Project Title]} \\[0.5em]
    \large Risk Analyst -- Frontier Quantitative Risk Analysis
}
\author{Andr\'es Gonz\'alez Ortega}
\date{\today}

\begin{document}
\maketitle

\begin{abstract}
    [One paragraph: what problem this project solves, which techniques it uses,
    and what the key results are.]
\end{abstract}

\tableofcontents
\newpage

% ──────────────────────────────────────────────────────────────────────────────
\section{Motivation and Problem Statement}
% Why does this problem matter? What real-world risk question does it address?
% Frame it for someone who understands probability but may not know the specific
% technique.

% ──────────────────────────────────────────────────────────────────────────────
\section{Mathematical Foundations}
% Full derivations. Don't skip steps. Number every important equation.

\subsection{Setup and Notation}
% Define the probability space, random variables, and key notation.

\subsection{Main Results}
% State and prove the core theorems/propositions.

\begin{keyresult}
    [State the main result concisely here.]
\end{keyresult}

% ──────────────────────────────────────────────────────────────────────────────
\section{Algorithm}
% Pseudocode for the computational method.

\begin{algorithm}[H]
\caption{[Algorithm Name]}
\begin{algorithmic}[1]
    \Require Input parameters
    \Ensure Output description
    \State [Step 1]
    \State [Step 2]
    \Return [Result]
\end{algorithmic}
\end{algorithm}

% ──────────────────────────────────────────────────────────────────────────────
\section{Connection to Code}
% Map equations to functions. Use \texttt{} for code references.

\begin{codelink}
    Equation~\eqref{} is implemented in \texttt{src/model.py:function\_name()},
    lines XX--YY.
\end{codelink}

% ──────────────────────────────────────────────────────────────────────────────
\section{Numerical Examples}
% Show specific calculations that match the code output.
% Include tables comparing analytic vs. computed results.

% ──────────────────────────────────────────────────────────────────────────────
\section{Extensions and Limitations}
% What assumptions were made? Where does the model break?
% What frontier research addresses these limitations?

% ──────────────────────────────────────────────────────────────────────────────
\section{References}
% Use biblatex or manual bibliography.

\end{document}
